\documentclass{svjour3}
\usepackage{graphicx}
\usepackage{amsmath}
\begin{document}
\title{Isolation and Characterization of Biosurfactant-Producing Enterobacter hormaechei subsp. hoffmannii MK-DU2 with Potential Application in Bioremediation of Crude Oil Contamination}
\author{Mina Kalantari and Fatemeh Salimi}
\institute{Department of Cellular and Molecular Biology, School of Biology, Damghan University, Damghan 36716-45667, Iran}
\maketitle
\begin{abstract}
Oil pollution has led to the growth of microbial communities producing biosurfactants that increase the bioavailability of petroleum hydrocarbons. This study aimed to isolate and characterize a new biosurfactant-producing strain as an alternative to toxic chemical surfactants.
\keywords{Biosurfactant, Emulsification, Hydrocarbon degradation, Soil bacteria}
\end{abstract}
% ===== MAIN DOCUMENT CONTENT =====

\section{Introduction}
The global surfactant industry, with annual production exceeding 7.5 million tons and generating over \$41 billion in revenue, faces increasing pressure to adopt sustainable alternatives to petrochemical-derived products \cite{ref1}. This shift is driven by environmental concerns associated with synthetic surfactants widely used across oil and gas, medicine, food production, agriculture, and personal care industries \cite{ref2}.
In response to these challenges, biosurfactants have emerged as promising eco-friendly alternatives due to their biodegradability, reduced toxicity, and functionality in extreme conditions, including temperature extremes, pH variations, and high salinity \cite{ref3}. These microbial compounds maintain remarkable activity even at low concentrations, enhancing their economic viability \cite{ref4}. Their amphiphilic structure—featuring a hydrophobic hydro-carbon tail and hydrophilic head composed of peptides, carbohydrates, amino acids, or phosphate carboxyl acids—enables exceptional surface-active properties at water-oil interfaces \cite{ref5}.
The environmental applications of biosurfactants have gained significant attention in recent years, particularly for bioremediation of hydrocarbon-contaminated environments. Recent studies have shown that lipopeptide biosurfactants show particular promise in bioremediation processes due to their ability to enhance the bioavailability of hydrophobic pollutants and facilitate microbial degradation of petroleum compounds \cite{ref6, ref7}. Of particular importance for environmental applications, studies have demonstrated that microbial biosurfactants can function effectively under extreme environmental conditions, including temperatures up to 70 $^\circ$C, pressures up to 6000 psia, salinities up to 15\% (w/v), and pH ranges of 4–10 \cite{ref8, ref9}.
Multi-component biosurfactants, containing combinations of lipid, peptide, and carbohydrate moieties, have shown particular promise for environmental applications.
Lipopeptide and protein-containing biosurfactants represent important categories of bio-surfactants, with lipopeptides being among the more abundant and widely studied, while protein-containing biosurfactants remain less explored but potentially valuable alternatives \cite{ref10}. Lipopeptides can be applied in diverse domains, including environmental applications due to their emulsification/de-emulsification capacity, and their qualification as viscosity reducers and hydrocarbon solubilizing agents \cite{ref11}. These compounds have been successfully used in enhancing oil recovery under extreme environmental conditions \cite{ref12}. Among the diverse microbial genera capable of biosurfactant production, Enterobacter species have emerged as promising candidates for environmental applications. Biosurfactants from Enterobacter species are beneficial for the bioremediation of hydrocarbons in the environment due to their solubilizing properties \cite{ref13}. Recent research has demonstrated that Enterobacter cloacae strains can serve as effective biosurfactant producers with potential for applications in the bioremediation of oil-contaminated sites \cite{ref14}.
This study addresses these knowledge gaps by isolating and characterizing Enterobacter hormaechei MK-DU2, a biosurfactant-producing strain from agricultural soil samples. Our objectives are to (1) evaluate the functional properties of the crude biosurfactant extract, including emulsification activity, surface tension reduction, and crude oil biodegradation enhancement; (2) provide preliminary chemical characterization of the biosurfactant composition using multiple analytical techniques; and (3) assess the potential of biosurfactant-producing Enterobacter hormaechei MK-DU2 for rapid bioremediation applications. This work contributes to sustainable biosurfactant research by demonstrating exceptional performance under challenging environmental conditions, which distinguishes it from previously reported Enterobacter-derived biosurfactants and supports its potential for practical environmental and industrial applications.
\section{Materials and Methods}
Sample Collection and Isolation of Hydrocarbon-Consuming Bacteria
A soil sample was obtained from a productive pistachio farm located in Raziabad, a village in the Ghahab Sarsar area of Amirabad district, Damghan City, Semnan Province, Iran. The sampling site’s coordinates were 36.11616585209723$^\circ$N and 54.21172885290508$^\circ$E.

The pistachio trees on the property were healthy, with no history of oil or hydrocarbon contamination. Sterilized glass containers, with a diameter of 10 cm and a height of 15 cm, were used to obtain the samples. Soil was extracted from a depth of 20 cm. After collection, the soil was carefully blended and filtered through a 2-mm sieve to ensure a uniform sample. To isolate culturable hydrocarbon-consuming bacteria, the vapor-phase method was used. The process involved placing small amounts of diluted samples onto specialized growth media. One hundred microliters of diluted sample (10$^{-5}$) were spread onto a modified mineral salt agar medium, containing MgSO$_4$.7H$_2$O (0.42 gL$^{-1}$), KCl (0.297 gL$^{-1}$), KH$_2$PO$_4$ (0.85 gL$^{-1}$), NaNO$_3$ (0.42 gL$^{-1}$), K$_2$HPO$_4$ (1.27 gL$^{-1}$), NaCl (0.1 gL$^{-1}$), and agar (20 gL$^{-1}$). Subsequently, diesel-soaked filter paper (Whatman No.1) was positioned on the inverted lid of each Petri dish, allowing hydrocarbons to reach the bacteria through evaporation. In this method, diesel served as the only source of carbon and energy for the bacteria. The samples were incubated at room temperature for approximately 1 week \cite{18, 19}.
\section{Screening of Selected Strains to Investigate Biosurfactant Production}
The bacteria were inoculated into Bushnell Haas Broth as a seeding medium, containing MgSO$_4$ (0.2 gL$^{-1}$), CaCl$_2$ (0.02 gL$^{-1}$), KH$_2$PO$_4$ (1 gL$^{-1}$), K$_2$HPO$_4$ (1 gL$^{-1}$), NH$_4$NO$_3$ (1 gL$^{-1}$), FeCl$_3$ (0.05 gL$^{-1}$), olive oil (10 mL), and pH 7. The inoculated media were incubated at 28 $^{\circ}$C for 24 h at 180 rpm. At the end of incubation, the seeding media were inoculated (1\% v/v) into fermentation media of the same composition. Finally, the flasks were incubated at 28 $^{\circ}$C for 96 h at 180 rpm. Upon completion of the incubation period, the inoculated fermentation broth was centrifuged at 13,000 rpm for 20 min at 4 $^{\circ}$C. The pellet containing the cells was discarded, and the cell-free supernatants (CFS) were utilized for the oil spreading and drop collapse assays \cite{20}.
\subsection{Oil Spreading Method}
An oil spreading assay was performed following the methodology outlined by Morikawa et al. \cite{21}. In this procedure, 40 mL of water was placed in a Petri dish, to which 50 $\mu$L of crude oil was added, resulting in a thin layer of oil on the water's surface. Subsequently, 40 $\mu$L of bacterial supernatant was introduced onto the oil, and the development of a clearing zone was observed. The ability of each isolate's fermentation liquid to disperse the oil and create a halo was assessed as a measure of biosurfactant production. The presence of the clearing zone was qualitatively linked to surfactant activity. Tween (1\% w/v) and an uninoculated culture medium served as positive and negative controls, respectively.
\subsection{Drop Collapse Method}
The assessment of biosurfactant production was conducted using the qualitative drop-collapse test, employing crude oil as the testing medium. In this procedure, 10 $\mu$L of crude oil were placed on the hydrophobic surface of parafilm and allowed to equilibrate for 24 h. Following this, 15 $\mu$L of cell-free culture broth was introduced onto the areas coated with crude oil, and the size of the resulting drop was examined 1 min later with the assistance of a magnifying glass. This approach facilitated a comparison of the spread of the drop against the positive control (Tween 1\% w/v) and the negative control (uninoculated culture medium). The shape of the droplet was observed after the addition of the supernatant, with rounded drops classified as negative, indicating an absence of biosurfactant production. Conversely, a flat drop was interpreted as a positive indication of biosurfactant production \cite{22}.
\subsection{Hemolysis Ability}
Bacterial isolates in pure culture were streaked onto freshly prepared blood agar and incubated at 37 $^{\circ}$C for 48 to 72 h. The results were documented according to the type of clear zone observed: $\alpha$-hemolysis was indicated by a greenish zone surrounding the colony, $\beta$-hemolysis was characterized by a clear, white zone around the colony, and $\gamma$-hemolysis was noted when there was no alteration in the medium surrounding the colony \cite{23}.
\subsection{Molecular Identification of the Biosurfactant-Producing Bacterium}
The strain exhibiting the highest efficiency in biosurfactant production was identified at the molecular level through 16S rRNA gene sequencing. Genomic DNA was isolated using a phenol-chloroform extraction method. Subsequently, the 16S rRNA gene was amplified via PCR using universal bacterial primers 27 F (AGA GTT TGA TCM TGG CTC AG) and 1492R (TAC GGY TAC CTT GTT ACG ACT T). The amplified product was purified and sequenced at a commercial facility (Pishgam, Iran) using Sanger sequencing. The resulting sequence was then compared against entries in the EzTaxon database for identification purposes. Afterwards, the sequence was submitted to GenBank, where it was assigned an accession number. Finally, a phylogenetic tree was constructed using MEGA-X (Molecular Evolutionary Genetics Analysis) software, applying the neighbor-joining method with 1000 bootstrap replicates \cite{24, 25}.
\subsection{Molecular Identification of the Biosurfactant-Producing Bacterium}
The strain exhibiting the highest efficiency in biosurfactant production was identified at the molecular level through 16S rRNA gene sequencing. Genomic DNA was isolated in duplicate from independent bacterial cultures using a phenol-chloroform extraction method to ensure reproducibility. Subsequently, the 16S rRNA gene was amplified via PCR in triplicate reactions using universal bacterial primers 27 F (AGA GTT TGA TCM TGG CTC AG) and 1492R (TAC GGY TAC CTT GTT ACG ACT T) \cite{26}. The amplified product was purified and subjected to bidirectional sequencing at a commercial facility (Pishgam, Iran) using Sanger sequencing with both forward and reverse primers to obtain high-quality, overlapping sequence data. The resulting consensus sequence was compared against entries in the EzTaxon database for identification purposes. The sequence was subsequently submitted to GenBank for accession number assignment. A phylogenetic tree was constructed using MEGA-X (Molecular Evolutionary Genetics Analysis) software, applying the neighbor-joining method with 1000 bootstrap replicates.
\subsection{Evaluating Emulsification Ability}

\section{Evaluating Emulsification Ability}
To assess the emulsification capability of the biosurfactant produced by the most effective bacterium, equal volumes of its CFS were combined with olive oil, sesame oil, diesel, or xylene. The resulting mixtures were vortexed for 2 min, after which the tubes were allowed to stand at room temperature for 1, 24, and 48 h. Emulsifying activity was measured at three distinct time points (EA1, EA24, and EA48) under varying conditions of acidity (pH 3, 6, and 8), salinity (NaCl concentrations of 1, 5, 10, and 20\% w/v), and monovalent salt presence (KCl concentrations of 1, 5, 10, and 20\% w/v). The evaluation of emulsion activity at different pH levels was conducted by adjusting the pH of the collected supernatant to the desired level using HCl (1M) and NaOH (1M), followed by vigorous mixing with an equal volume of the oils or solvents. The emulsifying activity was then measured. Similarly, emulsifying activity at varying concentrations of NaCl and KCl was assessed by adding the specified amounts of salts into the collected supernatant after 1, 24, and 48 h \cite{27}.
Emulsifying activity (EA) =
\[
\frac{\text{Height of the emulsion layer}}{\text{Overall height of the mixture}}
\]
$\times 100$
\subsection{Biosurfactant Extraction}
The pH of the CFS was adjusted to 2.0 by the addition of 6 M HCl and maintained at 4 $^\circ$C overnight. Subsequently, the mixture was centrifuged at 13,000 $\times$ g for 20 min at 4 $^\circ$C. The biosurfactant was then extracted from the supernatant three times using an equal volume of a 2:1 chloroform:methanol solution. After extraction, the organic phase was evaporated to obtain the crude biosurfactant. The resulting crude product underwent dialysis for 48 h at 4 $^\circ$C using dialysis tubing with a 12,000 Da molecular weight cut-off (Sigma) and subsequently lyophilized \cite{28}.
\subsection{Determining the Reduction of Surface Tension}
Measuring surface tension is an efficient and precise technique that uses a tensiometer, which relies on assessing the force required to detach a wire ring from a surface or from between two surfaces. This force correlates directly with surface tension. To evaluate the reduction of surface tension in the presence of a biosurfactant, 2 mg of the extracted biosurfactant was dissolved in 1 mL of Milli-Q water, and the surface tension of the resulting mixture was measured at 26 $^\circ$C using a digital K11–Kruss tensiometer with the Du Noüy ring method \cite{29, 30}.
\subsection{Hydrocarbon Degradation and Residual Crude Oil Determination}
Erlenmeyer flasks with a capacity of 250 mL were thoroughly cleaned and sterilized via autoclaving. Each flask was then filled with 100 mL of sterilized minimal salt medium, which comprised K$_2$HPO$_4$ (0.8 g/L), KH$_2$PO$_4$ (0.2 g/L), CaSO$_4$$\cdot$2H$_2$O (0.05 g/L), MgSO$_4$$\cdot$7H$_2$O (0.5 g/L), (NH$_4$)$_2$SO$_4$ (1 g/L), and FeSO$_4$ (0.01 g/L). The experimental design employed two distinct conditions to evaluate the hydrocarbon degradation capability of Enterobacter hormaechei MK-DU2 strain. The first condition consisted of inoculated flasks containing minimal salt medium, crude oil (5\% v/v), and freshly cultivated E. hormaechei MK-DU2 culture (2\% v/v) with an optical density (OD 625 nm) of 0.8, while the second condition served as an uninoculated control containing only minimal salt medium and crude oil without bacterial inoculation. The freshly cultivated bacterial culture was prepared by growing E. hormaechei MK-DU2 in nutrient broth for 18–24 h at 37 $^\circ$C with continuous shaking at 150 rpm until reaching the specified optical density. The pH of the flasks was adjusted to 7, and the experiments were conducted in a shaker incubator at 150 rpm for 6 days. This uninoculated control was essential to distinguish between abiotic degradation processes and biological degradation mediated by the inoculated strain, ensuring that observed hydrocarbon removal could be attributed to bacterial activity rather than physical or chemical processes. At 24-h intervals throughout the experimental period, residual crude oil was extracted from both inoculated and control flasks using an equal volume of hexane. The oil-hexane mixture was then carefully separated using a separating funnel and collected in pre-weighed tubes for gravimetric analysis. The amount of residual oil was determined by evaporating the hexane solvent at 60 $^\circ$C for 12 h, following established protocols \cite{31}.

with continuous shaking at 150 rpm until reaching the specified optical density. The pH of the flasks was adjusted to 7, and the experiments were conducted in a shaker incubator at 150 rpm for 6 days. This uninoculated control was essential to distinguish between abiotic degradation processes and biological degradation mediated by the inoculated strain, ensuring that observed hydrocarbon removal could be attributed to bacterial activity rather than physical or chemical processes. At 24-h intervals throughout the experimental period, residual crude oil was extracted from both inoculated and control flasks using an equal volume of hexane. The oil-hexane mixture was then carefully separated using a separating funnel and collected in pre-weighed tubes for gravimetric analysis. The amount of residual oil was determined by evaporating the hexane solvent at 60 $^\circ$C for 12 h, following established protocols \cite{31}.
The percentage of hydrocarbon degradation was calculated using the equation:
Hydrocarbon Degradation(\%) =
Weight of residual oil in uninoculated control $-$ Weight of residual oil in inoculated sample
Original weight of crude oil
$\times 100\%$
\subsection{Determining the Carbohydrate, Lipid, and Protein Nature of the Surfactant Using Chromatographic Analysis}
Chromatographic analysis was conducted to ascertain the carbohydrate, lipid, and protein composition of the extracted biosurfactant. Ten milligrams of the extracted biosurfactant was dissolved in 1 mL of chloroform, and approximately 40 µL of this solution was applied to a thin-layer chromatography (TLC) plate. The TLC was performed using a mobile phase composed of chloroform, methanol, acetic acid, and water in a ratio of 10:50:60:4. To identify proteins, sugars, and lipids, the plates were treated with 0.5\% ninhydrin solution (0.5 g of ninhydrin dissolved in 100 mL of ethanol) and anisaldehyde solution (prepared from acetic acid, sulfuric acid, and p-anisaldehyde in a ratio of 100:2:1) and exposed to iodine vapor in a saturated chamber, respectively. Following the staining process, the plates were incubated at 110 $^\circ$C for 10 min to facilitate the development of spots \cite{32}.
\subsection{Determining Functional Groups Using the Infrared Spectrometer (FTIR) Method}
To elucidate the chemical characteristics of biosurfactants through the identification of functional groups, Fourier transform infrared (FT-IR) spectroscopy (Equinox 55, Bruker, Germany) was employed. For this analysis, pellets with a diameter of 3 mm were prepared by combining KBr powder with lyophilized biosurfactants in a ratio of 100:1 (KBr:biosurfactants). A KBr pellet was used as the background reference. The FT-IR spectrum was then recorded in the range of 4000--400 cm$^{-1}$, with a scanning resolution of 4 cm$^{-1}$.
\subsection{Nuclear Magnetic Resonance Spectroscopy (1H-NMR) of the Biosurfactant}
The molecular characteristics of biosurfactants in solution were examined using nuclear magnetic resonance (NMR) spectroscopy, specifically focusing on hydrogen nuclei. The $^{1}$H NMR spectrum of the biosurfactant solution was obtained using a spectrometer. To this analysis, 5 mg of lyophilized biosurfactant was dissolved in 1 mL of deuterated chloroform (CDCl$_{3}$), and the chemical shifts ($\delta$) were reported in parts per million (ppm) at a frequency of 500 MHz.
\subsection{Thermal Stability of the Biosurfactant}
Thermogravimetric analysis (TGA) of biosurfactant was performed using a TGA instrument from Mettler Toledo, Switzerland. For this purpose, 35 mg of lyophilized biosurfactant was placed onto a platinum pan and heated from 25 to 600 $^\circ$C at a rate of 10 $^\circ$C min$^{-1}$ under a nitrogen atmosphere. The corresponding weight loss was determined.
\section{Result and Discussion}

\section{Result and Discussion}
\subsection{Qualitative Tests for Screening Biosurfactant-Producing Bacterial Isolates}
The oil spreading test, a reliable method for detecting biosurfactant production, revealed varying degrees of activity among the bacterial isolates tested (Fig.~1a). The CFS of MK-DU2 exhibited the highest activity with a clear zone diameter of 21~mm, followed by the CFSs of MK-DU1 and MK-DU3, with diameters of 19~mm and 17~mm, respectively. MK-DU4 demonstrated multiple smaller clear zones, indicating lower biosurfactant production. Although limited comparative data exist in the literature, these findings align with previous studies on \textit{Enterobacter cloacae} B14 and \textit{Enterobacter quasihormaechei} \cite{14, 33}. While these earlier studies confirmed biosurfactant production capability in related \textit{Enterobacter} species, they did not quantify displacement zone diameters, making our measurements a valuable contribution to the existing knowledge base of \textit{Enterobacter} biosurfactant characterization.
The drop collapse assay provided complementary evidence of biosurfactant production. CFSs from all four bacterial isolates produced significantly flattened and spread droplets compared to the negative control (Fig.~1b).
Among these isolates, only MK-DU2 exhibited a substantial zone of partial hemolysis (alpha-hemolysis) in the blood agar test (Fig.~1c). The varying degrees of hemolysis among the isolates suggest differences in either the type or quantity of biosurfactants produced. It is important to note that while hemolysis is a useful initial indicator, it is not definitive proof of biosurfactant production. Some hemolytic factors may not be biosurfactants, and some biosurfactants may not cause hemolysis \cite{34}. Therefore, these results should be considered preliminary, and further biochemical tests and structural analyses are necessary to confirm and characterize the biosurfactants produced by these isolates.
\subsection{Molecular Identification of the Biosurfactant-Producing Isolate MK-DU2}
Molecular characterization through replicated 16S rRNA gene sequencing consistently determined that isolate MK-DU2 (GenBank accession number PQ283906) shares 99.45\% sequence identity with \textit{Enterobacter hormaechei} subsp. \textit{hoffmannii} across all technical replicates. Biosurfactant production has been documented in several \textit{Enterobacter} species, with \textit{E. hormaechei} strains specifically known to synthesize rhamnolipids (Rha-Rha-C10-C12) \cite{35} and anionic lipoprotein-type biosurfactants \cite{36}. This reproducible phylogenetic placement strongly suggests that our MK-DU2 isolate likely produces similar glycolipid or lipoprotein biosurfactant compounds. Our findings expand the current understanding of biosurfactant biosynthetic capabilities within the \textit{Enterobacter} genus, particularly for the subspecies \textit{E. hormaechei} subsp. \textit{hoffmannii}, which has been less extensively characterized for this biotechnological application.
\subsection{Emulsification Ability of Enterobacter hormaechei MK-DU2 in Harsh Conditions}
To provide a more comprehensive characterization, the emulsification ability of the biosurfactant derived from \textit{Enterobacter hormaechei} MK-DU2 was investigated with four hydrocarbons—olive oil, sesame oil, diesel, and xylene—across different pH levels and in the presence of varying concentrations of salinity and monovalent salts.
\subsubsection{Emulsifying Activity Across Diverse pH Conditions}
The biosurfactant extract from \textit{E. hormaechei} MK-DU2 demonstrated significant emulsification activity across various pH levels with four different hydrocarbons (Fig.~2). Emulsification patterns varied distinctly depending on both pH conditions and hydrocarbon substrate. For sesame oil and diesel, optimal activity occurred at pH 6 and 8, indicating preferential performance under neutral to slightly alkaline conditions. In contrast, xylene and olive oil exhibited enhanced emulsification at acidic pH, with olive oil maintaining consistently high indices across all tested pH ranges. Also, the amount of emulsion in the absence of biosurfactant was represented in Table~S1.
\begin{figure}[htbp]
\centering
\includegraphics[width=0.8\linewidth]{images/figure2.png}
\caption{Effect of pH on emulsifying activity of biosurfactant containing cell-free supernatant of \textit{Enterobacter hormaechei} MK-DU2 on four types of oils after 1, 24, and 48 h}
\label{fig:2}
\end{figure}

The biosurfactant-containing cell-free supernatant of \textit{Enterobacter hormaechei} MK-DU2 consistently showed high indices across all tested pH ranges. Also, the amount of emulsion in the absence of biosurfactant was represented in Table S1.
These results both confirm and challenge established patterns in \textit{Enterobacter} biosurfactant research. The neutral-to-alkaline preference observed for sesame oil and diesel emulsification aligns with the broader literature. For example, biosurfactant production from \textit{Enterobacter cloacae} AYF1 was optimal at pH 7.5, with stability tests indicating strong emulsification activity around neutral pH (pH 7) and at temperatures near 40--45 $^{\circ}$C \cite{37}. This suggests that many \textit{Enterobacter}-derived biosurfactants are stable and active under neutral to slightly alkaline conditions. Similarly, research on \textit{Enterobacter cloacae} B14 reported improved biosurfactant production with stability and antimicrobial activity in neutral to alkaline pH \cite{37}. In addition, a study characterizing \textit{Enterobacter cloacae} BAGM01 from the Gulf of Mexico described a thermostable and alkaline-tolerant rhamnolipid biosurfactant \cite{16}.
Notably, while most reported \textit{Enterobacter}-derived biosurfactants show optimal activity at neutral to alkaline pH, the biosurfactant-containing CFS from MK-DU2 exhibited exceptional emulsification of xylene and olive oil under acidic conditions. This acid stability represents a distinctive property for biosurfactants produced by this \textit{Enterobacter} strain. These findings contrast with those of Lovaglio et al. (2011), who reported the breakdown of rhamnolipid emulsions after prolonged exposure to acidic environments \cite{38}, but align with the acid stability of a lipopeptide biosurfactant from \textit{Enterobacter cloacae} UCP 1597 \cite{17}.
This pH-dependent behavior can be explained by several interconnected mechanisms. pH affects the ionization state of biosurfactant functional groups, directly influencing their charge distribution and solubility \cite{39}. These changes alter micelle formation and critical micelle concentration \cite{41}, with extreme pH values potentially leading to precipitation and reduced biosurfactant availability for surface activity \cite{40}. Additionally, pH-induced structural changes can affect the conformation of biosurfactant molecules---particularly in protein-based biosurfactants---thereby significantly impacting their surface-active properties \cite{41}.
\begin{figure}[htbp]
\centering
\includegraphics[width=0.8\linewidth]{images/figure3.png}
\caption{Effect of NaCl in various concentrations on emulsifying activity of biosurfactant containing cell-free supernatant of \textit{Enterobacter hormaechei} MK-DU2 on four types of oils after 1, 24, and 48 h}
\label{fig:3}
\end{figure}
The interfacial behavior of biosurfactants is inherently pH-dependent, influencing their adsorption at interfaces due to changes in molecular and interfacial charges \cite{42}. Emulsion stability is likewise affected by pH through alterations in electrostatic repulsions between droplets and changes in interfacial film properties \cite{43}. Furthermore, pH modifies the interaction of biosurfactants with other molecules in the system, thereby impacting their surface activity \cite{44}.
\subsection{Emulsifying Activity in the Presence of Salinity}
The biosurfactant-containing supernatant of \textit{Enterobacter hormaechei} MK-DU2 maintained significant emulsification activity at NaCl concentrations up to 5\%, with some activity persisting even at 20\% NaCl (Fig. 3). This level of salt tolerance exceeds that reported for most \textit{Enterobacter}-derived biosurfactants. Also, the amount of emulsion in the absence of biosurfactant was represented in Table S2. For comparison, \textit{E. cloacae} BAGM01 produced a biosurfactant that remained stable at 3\% NaCl, whereas our extract maintained considerable activity at twice this concentration \cite{16}. The emulsions remained stable at lower salt concentrations (1--5\% NaCl) over 48 h, indicating the production of robust biosurfactants suitable for industrial and environmental applications such as bioremediation in coastal areas. At higher salt concentrations (10--20\% NaCl), emulsion stability decreased more noticeably over time, with olive oil consistently demonstrating the highest activity.
\begin{figure}[htbp]
\centering
\includegraphics[width=0.8\linewidth]{images/figure4.png}
\caption{Effect of KCl in various concentrations on emulsifying activity of biosurfactant containing cell-free supernatant of \textit{Enterobacter hormaechei} MK-DU2 on types of oils after 1, 24, and 48 h}
\label{fig:4}
\end{figure}
The ability to maintain activity in 6\% NaCl is comparable to that of the biosurfactant from \textit{Virgibacillus salarius} KSA-T under similar conditions \cite{45}, while retention of activity even at 10\% NaCl indicates the production of moderately halotolerant biosurfactants. The observed decrease in biosurfactant activity with increasing salinity can be attributed to several mechanisms: the salting-out effect, which reduces biosurfactant solubility \cite{46}, electrostatic screening by salt ions diminishing repulsion between charged biosurfactant molecules \cite{42}, alterations in the critical micelle concentration \cite{47}, modifications to the biosurfactant's molecular structure \cite{48}, competition with biosurfactants for water molecules \cite{49}, and changes in interfacial properties \cite{50}.

\section{Emulsifying Activity in the Presence of Monovalent Salt}
The biosurfactant-containing supernatant also demonstrated remarkable tolerance to potassium chloride (Fig. 4). Although emulsification activity generally decreased with increasing KCl concentration, significant activity was maintained at 1\% and 5\% KCl for most oils. Also, the amount of emulsion in the absence of biosurfactant was represented in Table S3. Previous studies on biosurfactants from other bacteria, such as Lactobacillus strains, have reported varying effects of salts on emulsification \cite{51}.
At 1\% KCl, all oils demonstrated high initial emulsifying activity, which remained relatively stable over 48 h. For example, olive oil maintained an emulsion stability (ES) of 76.92\% after 48 h at 1\% KCl. Olive oil emulsions exhibited the highest stability across all KCl concentrations during the first hour, with an EA of 55\% at 20\% KCl. However, these emulsions showed a decline in the presence of 10\% and 20\% KCl after 24 and 48 h, with the EA dropping to 12.72\% at 20\% KCl after 48 h. Diesel emulsions displayed moderate stability, with some separation observed at higher KCl concentrations. Xylene emulsions demonstrated the least stability, with rapid separation occurring at all salt concentrations.
Interestingly, xylene showed an unexpected spike in EA (reaching 66.6\%) at 5\% KCl after 1 h, which decreased rapidly over time. This anomaly could be due to a temporary stabilization effect of KCl at this specific concentration, possibly related to changes in interfacial tension or electrostatic interactions. Further investigation would be needed to fully explain this phenomenon.
\begin{figure}[htbp]
\centering
\includegraphics[width=0.8\linewidth]{images/figure1.png}
\caption{Enhancement of petroleum hydrocarbon biodegradation by biosurfactant-producing Enterobacter hormaechei MK-DU2 during 6 days using a gravimetric method}
\label{fig:label}
\end{figure}
The observed salt-specific effects may have implications for applications requiring the targeting of specific hydrocarbons in mixed contamination scenarios. The enhanced activity with specific hydrocarbon-salt combinations represents a noteworthy finding.
\subsection{Surface Tension Reduction by Enterobacter hormaechei MK-DU2 Biosurfactant Extract}
The crude biosurfactant extract derived from Enterobacter hormaechei MK-DU2 effectively reduced the surface tension of water from 69.7 to 44.4 mN/m, achieving a significant 36.3\% decrease. Although this level of surface tension reduction is comparable to previously reported biosurfactants—for instance, Legawiec et al. (2023) demonstrated rhamnolipids reducing surface tension to approximately 32–33 mN/m at critical micelle concentration (CMC, 26–90 mg/L) \cite{52}, and Santos et al. documented Streptomyces-derived biosurfactants lowering it to ~33.76 mN/m at CMC (~10 mg/mL) \cite{53}—what distinctly sets our biosurfactant apart is its remarkable stability and sustained activity under extreme environmental conditions. Specifically, this biosurfactant maintains its surface tension reducing capability across a broad range of pH values and in the presence of high salt concentrations, conditions which frequently diminish the efficacy of most biosurfactants. This unique combination of moderate, yet consistent, surface activity and robust functional stability under harsh acidic and saline environments provides a significant advantage for practical applications, particularly in complex environmental and industrial settings where variable and extreme conditions are common.
\subsection{Enhanced Petroleum Hydrocarbon Biodegradation by Enterobacter hormaechei MK-DU2: Combined Biosurfactant Production and Direct Metabolism}
The hydrocarbon degradation capability of Enterobacter hormaechei MK-DU2 was comprehensively evaluated through direct strain inoculation experiments conducted over 6 days, as illustrated in Fig. 5. The experimental design compared petroleum hydrocarbon degradation between inoculated flasks containing E. hormaechei MK-DU2 and uninoculated control flasks to distinguish biological degradation from any potential abiotic processes. Inoculation with E. hormaechei MK-DU2 resulted in remarkable biodegradation performance, achieving 53.34\% degradation within the first 24 h and progressing to an impressive 99.23\% degradation by day 6. This exceptional degradation efficiency demonstrates the outstanding biodegradation capability of the MK-DU2 strain and establishes its potential as a highly effective bioremediation agent.
The superior biodegradation performance observed with E. hormaechei MK-DU2 likely results from the strain’s ability to both metabolize petroleum hydrocarbons and produce biosurfactants during growth, as previously characterized in our biosurfactant production studies. The enhanced degradation rates suggest that the strain may utilize multiple complementary mechanisms during the biodegradation process. Based on literature reports for similar Enterobacter species, the degradation process potentially involves direct metabolic utilization of hydrocarbon compounds as carbon and energy sources, coupled with the production of biosurfactants that could enhance substrate bioavailability. These biologically produced surfactants potentially enhance the biodegradation process by reducing surface tension at the oil-water interface, forming micelles that may solubilize hydrophobic petroleum compounds, and increasing the effective surface area available for microbial activity. Visual observations during the experiment supported possible biosurfactant activity, including apparent reduction in oil-water surface tension and formation of stable emulsions in inoculated flasks.

The superior biodegradation performance observed with E. hormaechei MK-DU2 likely results from the strain’s ability to both metabolize petroleum hydrocarbons and produce biosurfactants during growth, as previously characterized in our biosurfactant production studies. The enhanced degradation rates suggest that the strain may utilize multiple complementary mechanisms during the biodegradation process. Based on literature reports for similar Enterobacter species, the degradation process potentially involves direct metabolic utilization of hydrocarbon compounds as carbon and energy sources, coupled with the production of biosurfactants that could enhance substrate bioavailability. These biologically produced surfactants potentially enhance the biodegradation process by reducing surface tension at the oil-water interface, forming micelles that may solubilize hydrophobic petroleum compounds, and increasing the effective surface area available for microbial activity. Visual observations during the experiment supported possible biosurfactant activity, including apparent reduction in oil-water surface tension and formation of stable emulsions in inoculated flasks.
The biodegradation efficiency achieved by E. hormaechei MK-DU2 substantially exceeds performance levels reported in the literature for comparable hydrocarbon-degrading microorganisms. While Pseudomonas aeruginosa PG1 achieved 81.8\% crude oil degradation over 5 weeks \cite{31}, our strain accomplished 99.23\% degradation in merely 6 days, representing an approximately sixfold improvement in degradation rate. Even when compared to other Enterobacter hormaechei strains that have demonstrated 88.2\% degradation efficiency over 35 days \cite{54}, the MK-DU2 strain shows remarkable superiority in both total degradation percentage and time efficiency. This exceptional performance likely results from the optimal integration of high intrinsic hydrocarbon-degrading enzyme activity with efficient biosurfactant production capability, facilitated by growth conditions that promote both metabolic and biosurfactant synthesis pathways. The statistical analysis of triplicate experiments, represented by error bars in Fig.~5, demonstrates significant differences between inoculated and control conditions (p < 0.05), confirming the reliability and reproducibility of the observed biodegradation enhancement. These results establish E. hormaechei MK-DU2 as an exceptionally effective single-strain solution for petroleum hydrocarbon bioremediation, combining rapid degradation kinetics with near-complete hydrocarbon removal efficiency, making it particularly suitable for emergency response applications where both speed and thoroughness of remediation are critical factors.
\begin{figure}[htbp]
\centering
\includegraphics[width=0.8\linewidth]{images/figure6.png}
\caption{Thin layer chromatography (TLC) characterization of crude biosurfactant extract from Enterobacter hormaechei MK-DU2.}
\label{fig:6}
\end{figure}
The crude biosurfactant extract was separated on silica gel TLC plates and visualized using specific staining reagents to identify different molecular components. a TLC plate stained with iodine vapor for lipid detection, showing distinct yellow-brown spots (highlighted by dashed rectangles) in the upper and lower regions indicating the presence of lipid moieties with different polarities and Rf values. b TLC plate stained with ninhydrin reagent for protein or peptide detection, displaying purple-pink colored spots at various positions on the plate, confirming the presence of amino acid or peptide components within the biosurfactant structure. c TLC plate stained with p-anisaldehyde reagent for carbohydrate detection, revealing dark-colored spots that indicate the presence of sugar moieties in the biosurfactant compound. The multiple positive reactions across all three detection methods demonstrate that the biosurfactant produced by MK-DU2 is a complex amphiphilic molecule containing lipid, peptide, and carbohydrate components
\subsection{Chemical Characterization of the Crude Biosurfactant Extract from Enterobacter hormaechei MK-DU2}
\subsubsection{Thin Layer Chromatography Analysis}
TLC analysis of the crude biosurfactant extract revealed the presence of multiple components, including lipid, peptide, and carbohydrate moieties (Fig.~6). Two distinct lipid spots were observed on the TLC plate: the upper spot likely represents more non-polar lipids, while the lower spot may correspond to slightly more polar lipid components. This indicates that the crude extract contains a complex mixture of lipids with varying polarities, consistent with the amphiphilic nature required for surface-active properties.
Multiple peptide bands were also visible, suggesting that the extract contains a mixture of peptides or small proteins. Additionally, a spot that reacted positively with p-anisaldehyde indicated the presence of carbohydrate components. Overall, these results suggest that the crude biosurfactant extract contains multiple components, potentially including lipids, peptides, and carbohydrates. Further purification and structural analysis are necessary to determine whether these components correspond to structures similar to or distinct from the rhamnolipids and lipoproteins previously reported from Enterobacter species \cite{35, 36}. It should be noted that further purification and structural analysis would be necessary to definitively characterize the specific components present. Notably, similar multi-component biosurfactant mixtures have been reported from Lactiplantibacillus plantarum OL5 \cite{55}, Aeromonas hydrophila RP1 \cite{56}, and Lactobacillus pentosus NCIM 2912 \cite{57}.
\begin{figure}[htbp]
\centering
\includegraphics[width=0.8\linewidth]{images/figure7.png}
\caption{FTIR spectrum of crude biosurfactant extract from Enterobacter hormaechei MK-DU2}
\label{fig:7}
\end{figure}
\subsubsection{FT-IR Analysis of Functional Groups in Enterobacter hormaechei MK-DU2-Derived Biosurfactant}
The FT-IR spectrum reveals absorption bands that provide preliminary information about functional groups present in our crude extract (Fig.~7). The peaks at 3389.12 cm$^{-1}$ may indicate the presence of hydroxyl (-OH) and/or amine (-NH) groups, while peaks at 3009.48 cm$^{-1}$ and 2925.67 cm$^{-1}$ could correspond to aliphatic C--H bonds in methyl and methylene groups. Peaks at 1745.32 cm$^{-1}$ and 1459.05 cm$^{-1}$ may be attributed to carbonyl and ester groups. Peaks at 1142.07 cm$^{-1}$ and 1054.18 cm$^{-1}$ are likely related to C-O-C stretching vibrations, which are characteristic of carbohydrates and other organic compounds \cite{57}. Additional peaks at 709.79 cm$^{-1}$, 1397.93 cm$^{-1}$, and 2854.34 cm$^{-1}$ further support the presence of long aliphatic fatty acid chains \cite{57}. The peak at 1236.32 cm$^{-1}$ could be attributed to P=O stretching vibration \cite{5}. Absorption peaks in the range of 1500--1700 cm$^{-1}$ are generally assigned to asymmetric stretching vibrations of C=O and carboxyl groups (-COO-) in fatty acids and amino acids, or to N-H groups in proteins.
This spectral profile provides preliminary evidence for the presence of various functional groups in the crude extract. While rhamnolipids from E. hormaechei B4 exhibited peaks primarily associated with glycolipid structures \cite{35}, the spectrum of our extract indicates a more complex composition with peaks corresponding to a wider range of organic functional groups. Additional purification and structural analysis will be necessary to definitively characterize the specific biosurfactant structures present.

\section{1H NMR Analysis of Crude Biosurfactant Produced by Enterobacter hormaechei MK-DU2}
The $1$H NMR spectrum of the crude biosurfactant extract (Fig.~\ref{fig:8}) provided further preliminary insight into its composition. Signals corresponding to aliphatic chains (peaks below $1.0$ ppm and between $1$ and $1.4$ ppm), possible unsaturated bonds (signals between $1.8$ and $2.2$ ppm), and hydroxyl groups (signals between $4$ and $4.2$ ppm) were observed, suggesting the presence of various organic moieties \cite{ref58}. Notably, signals at $5.341$ and $5.359$ ppm may indicate unsaturated bonds (possibly $-CH=CH-$ groups), commonly found in fatty acid chains of biosurfactants. A peak at $4.152$ ppm suggests the presence of carbohydrate components, potentially part of glycolipid structures, although further analysis is needed to confirm this association. The characteristic solvent peak (CDCl$_3$) was observed at $7.26$ ppm.
These spectral features are consistent with those reported for biosurfactants produced by other bacteria. For instance, Pandey et al. (2022) and Sharma (2022) identified similar peaks in the aliphatic and carbohydrate regions for biosurfactants from Aeromonas hydrophila and Lactobacillus pentosus NCIM 2912 \cite{ref56, ref57}. It is important to note that this NMR analysis was conducted on the crude extract without extensive purification. While these data provide valuable preliminary information regarding the types of functional groups present, they are insufficient to definitively determine specific molecular structures. Further purification and additional analytical techniques—such as mass spectrometry, $2$D NMR analyses, and isolation of individual components—will be necessary to fully characterize the specific biosurfactant molecules produced by E. hormaechei MK-DU2.
\begin{figure}[htbp]
\centering
\includegraphics[width=0.8\linewidth]{images/figure8.png}
\caption{1H nuclear magnetic resonance spectrum of crude biosurfactant extract produced by Enterobacter hormaechei MK-DU2}
\label{fig:8}
\end{figure}
\section{Thermal Stability Analysis of Crude Biosurfactant Produced by Enterobacter hormaechei MK-DU2}
TGA analysis demonstrated that the biosurfactant extract maintained stability up to $100$ $^\circ$C, after which a two-stage decomposition process was observed (Fig.~\ref{fig:9}). In the first stage ($100$–$300$ $^\circ$C), a substantial weight loss of $60\%$ occurred. This loss can be attributed to the evaporation of residual water and solvents, as well as the decomposition of heat-sensitive components such as lipids, sugar moieties, and proteins. The observed pattern suggests the presence of water-soluble constituents and potentially some bound water molecules within the biosurfactant. The major weight loss in this stage also indicates a significant lipid fraction and the decomposition of proteins and polysaccharide side chains \cite{ref59}.
\begin{figure}[htbp]
\centering
\includegraphics[width=0.8\linewidth]{images/figure9.png}
\caption{Thermogravimetric analysis of crude biosurfactant extract produced by Enterobacter hormaechei MK-DU2}
\label{fig:9}
\end{figure}
In the second stage ($300$–$500$ $^\circ$C), the remaining $40\%$ of the biosurfactant decomposed, likely representing the degradation of more thermally stable components within the crude extract. A similar TGA profile was reported by Gharaie et al. for the glycolipopeptide biosurfactant from Bacillus pumilus SG50, though notable differences were observed \cite{ref58}. While Gharaie et al. noted maximum weight loss at $280$ $^\circ$C due to rhamnose decomposition, our sample exhibited a more gradual weight loss between $220$ and $500$ $^\circ$C in the second stage, indicating differences in structural composition or thermal stability compared to rhamnolipid biosurfactant. The observed thermal stability of the crude biosurfactant extract, combined with its acid and salt tolerance, suggests a promising combination of properties for potential applications that require robust performance under challenging conditions.
\section{Conclusion}
In this study, we isolated and characterized a biosurfactant-producing strain, Enterobacter hormaechei MK-DU2, from soil samples. Our investigations revealed that the biosurfactant preparation from this strain exhibits remarkable functional properties, including strong emulsification activity and stability across a wide range of pH levels and salt concentrations, particularly in emulsifying olive oil. The crude biosurfactant extract achieved a substantial $36.3\%$ reduction in surface tension and significantly enhanced crude oil biodegradation, facilitating $99.23\%$ degradation within just $6$ days—a rate considerably faster than previously reported for similar applications.
Preliminary chemical analyses of the crude extract, including TLC, FTIR, and $1$H-NMR, indicated a complex composition comprising lipid, peptide, and carbohydrate components. The extract also demonstrated notable thermal stability up to $100$ $^\circ$C, followed by a two-stage decomposition process. However, further purification and advanced structural analyses are necessary to definitively characterize the specific biosurfactant molecules produced by this strain. Our findings contribute to the growing body of knowledge on biosurfactants produced by Enterobacter species, particularly highlighting the exceptional performance of the E. hormaechei MK-DU2 biosurfactant under high salinity and acidic pH conditions. This combination of functional properties suggests promising potential for applications in challenging environmental settings where conventional biosurfactants might be less effective.

Preliminary chemical analyses of the crude extract, including TLC, FTIR, and $1$H-NMR, indicated a complex composition comprising lipid, peptide, and carbohydrate components. The extract also demonstrated notable thermal stability up to $100$ $^\circ$C, followed by a two-stage decomposition process. However, further purification and advanced structural analyses are necessary to definitively characterize the specific biosurfactant molecules produced by this strain. Our findings contribute to the growing body of knowledge on biosurfactants produced by \textit{Enterobacter} species, particularly highlighting the exceptional performance of the \textit{E. hormaechei} MK-DU2 biosurfactant under high salinity and acidic pH conditions. This combination of functional properties suggests promising potential for applications in challenging environmental settings where conventional biosurfactants might be less effective.
Supplementary Information The online version contains supplementary material.
Acknowledgements The authors would like to thank Damghan University for its financial support of this research.
Author Contributions M.K. did the methodology and investigation, and F.S. did the conceptualization, supervision, and writing—review and editing.
Data Availability Data will be made available on request.
Declarations 
Ethics Approval Not applicable.
Consent to Participate The authors have contributed substantially to the manuscript and approved the final submission.
Consent for Publication The authors have approved the publication rules.
Competing Interests The authors declare no competing interests.

% ===== REFERENCES =====

\begin{thebibliography}{99}

\bibitem{ref1} Marti, M., Colonna, W., Patra, P., Zhang, H., Green, C., Reznik, G., et al. (2014). Production and char-acterization of microbial biosurfactants for potential use in oil-spill remediation. Enzyme and Micro-bial Technology, 55, 31–39. https:// doi. org/ 10. 1016/j. enzmi ctec. 2013. 09. 012

\bibitem{ref2} Dias, M. A. M., & Nitschke, M. (2023). Bacterial-derived surfactants: An update on general aspects and forthcoming applications. Brazilian Journal of Microbiology, 54(1), 103–123. https:// doi. org/ 10. 1007/ s42770- 023- 00905-7

\bibitem{ref3} Sundaram, T., Govindarajan, R. K., Vinayagam, S., Krishnan, V., Nagarajan, S., Gnanasekaran, G. R., et al. (2024). Advancements in biosurfactant production using agro-industrial waste for industrial and environmental applications. Frontiers in Microbiology, 15, 1357302. https:// doi. org/ 10. 3389/ fmicb. 2024. 13573 02

\bibitem{ref4} Nikolova, C., & Gutierrez, T. (2021). Biosurfactants and their applications in the oil and gas industry: Current state of knowledge and future perspectives. Frontiers in Bioengineering and Biotechnology, 9, Article 626639. https:// doi. org/ 10. 3389/ fbioe. 2021. 626639 7428 Applied Biochemistry and Biotechnology (2025) 197:7409–7430

\bibitem{ref5} Asgher, M., Arshad, S., Qamar, S. A., & Khalid, N. (2020). Improved biosurfactant production from Aspergillus niger through chemical mutagenesis: Characterization and RSM optimization. SN Applied Sciences, 2, 1–11. https:// doi. org/ 10. 1007/ s42452- 020- 2637-0

\bibitem{ref6} Santos, D. K. F., Rufino, R. D., Luna, J. M., Santos, V. A., & Sarubbo, L. A. (2016). Biosurfactants: Multifunctional biomolecules of the 21st century. International Journal of Molecular Sciences, 17(3), 401. https:// doi. org/ 10. 3390/ ijms1 70304 01

\bibitem{ref7} Silva, Rd. C. F., Almeida, D. G., Rufino, R. D., Luna, J. M., Santos, V. A., & Sarubbo, L. A. (2014). Applications of biosurfactants in the petroleum industry and the remediation of oil spills. International Journal of Molecular Sciences, 15(7), 12523–12542. https:// doi. org/ 10. 3390/ ijms1 50712 523

\bibitem{ref8} Bezerra, K. G. O., Rufino, R. D., Luna, J. M., & Sarubbo, L. A. (2018). Saponins and microbial bio-surfactants: Potential raw materials for the formulation of cosmetics. Biotechnology Progress, 34(6), 1482–1493. https:// doi. org/ 10. 1002/ btpr. 2682

\bibitem{ref9} Darvishi, P., Ayatollahi, S., Mowla, D., & Niazi, A. (2011). Biosurfactant production under extreme environmental conditions by an efficient microbial consortium, ERCPPI-2. Colloids and Surfaces B: Biointerfaces, 84(2), 292–300. https:// doi. org/ 10. 1016/j. colsu rfb. 2011. 01. 011

\bibitem{ref10} Bezerra, K. G. O., Rufino, R. D., Luna, J. M., & Sarubbo, L. A. (2018). Saponins and microbial bio-surfactants: Potential raw materials for the formulation of cosmetics. Biotechnology Progress, 34(6), 1482–1493. https://doi.org/10.1002/btpr.2682

\bibitem{ref11} Darvishi, P., Ayatollahi, S., Mowla, D., & Niazi, A. (2011). Biosurfactant production under extreme environmental conditions by an efficient microbial consortium, ERCPPI-2. Colloids and Surfaces B: Biointerfaces, 84(2), 292–300. https://doi.org/10.1016/j.colsurfb.2011.01.011

\bibitem{ref12} Ndlovu, T., Rautenbach, M., Vosloo, J. A., Khan, S., & Khan, W. (2017). Characterisation and antimicrobial activity of biosurfactant extracts produced by Bacillus amyloliquefaciens and Pseudomonas aeruginosa isolated from a wastewater treatment plant. AMB Express, 7(1), 108. https://doi.org/10.1186/s13568-017-0363-8

\bibitem{ref13} Singh, R., Glick, B. R., & Rathore, D. (2018). Biosurfactants as a biological tool to increase micro-nutrient availability in soil: A review. Pedosphere, 28(2), 170–189. https://doi.org/10.1016/S1002-0160(18)60020-7

\bibitem{ref14} Zhang, J., Xue, Q., Gao, H., Lai, H., & Wang, P. (2016). Production of lipopeptide biosurfactants by Bacillus atrophaeus 5–2a and their potential use in microbial enhanced oil recovery. Microbial Cell Factories, 15(1), 168. https://doi.org/10.1186/s12934-016-0587-3

\bibitem{ref15} Gharaei-Fathabad, E. (2011). Biosurfactants in pharmaceutical industry: A mini-review. American Journal of Drug Discovery and Development, 1(1), 58–69.

\bibitem{ref16} Ekprasert, J., Laopila, K., & Kanakai, S. (2019). Biosurfactant production by a newly isolated Enterobacter cloacae B14 capable of utilizing spent engine oil. Polish Journal of Environmental Studies, 28(4), 2603–2610. https://doi.org/10.15244/pjoes/92120

\bibitem{ref17} Jui, A. H., Bhuiyan, M. N. I., Bhowmik, B., Khatun, N., Chowdhury, A., Bhuiyan, R. H., et al. (2024). Exploration and characterization of a newly isolated bacterium, Enterobacter quasihormaechei strain BDIFST24001, capable of producing rhamnolipids biosurfactant for oil remediation. Access Microbiology, 000830, Article v1. https://doi.org/10.1099/acmi.0.000830

\bibitem{ref18} Curiel-Maciel, N. F., Martínez-Morales, F., Licea-Navarro, A. F., Bertrand, B., Aguilar-Guadarrama, A. B., Rosas-Galván, N. S., et al. (2021). Characterization of Enterobacter cloacae BAGM01 producing a thermostable and alkaline-tolerant rhamnolipid biosurfactant from the Gulf of Mexico. Marine Biotechnology, 23, 106–126. https://doi.org/10.1007/s10126-020-10033-4

\bibitem{ref19} de Souza Fonseca, T. C., Rodríguez, D. M., de Souza Mendonça, R., de Souza, A. F., Costa, L. O., & de Campos-Takaki, G. M. (2022). Multifunctional, stable and low-cost lipopeptide biosurfactant produced by Enterobacter cloacae UCP 1597. Research, Society and Development, 11(15), Article e41111536353. https://doi.org/10.33448/rsd-v11i15.36353

\bibitem{ref20} Curiel-Maciel, N. F., Martínez-Morales, F., Licea-Navarro, A. F., Bertrand, B., Aguilar-Guadar-rama, A. B., Rosas-Galván, N. S., et al. (2021). Characterization of Enterobacter cloacae BAGM01 producing a thermostable and alkaline-tolerant rhamnolipid biosurfactant from the Gulf of Mexico. Marine Biotechnology, 23, 106–126. https:// doi. org/ 10. 1007/ s10126- 020- 10033-4

\bibitem{ref21} de Souza Fonseca, T. C., Rodríguez, D. M., de Souza Mendonça, R., de Souza, A. F., Costa, L. O., & de Campos-Takaki, G. M. (2022). Multifunctional, stable and low-cost lipopeptide biosurfactant produced by Enterobacter cloacae UCP 1597. Research, Society and Development, 11(15), Article e41111536353. https:// doi. org/ 10. 33448/ rsd- v11i15. 36353

\bibitem{ref22} Soghandi, B., & Salimi, F. (2024). Study on amendment of rapeseed meal, soybean meal, and NPK fertilizer as biostimulants in bioremediation of diesel-contaminated soil by autochthonous microor-ganisms. Soil and Sediment Contamination, 33(4), 430–450.

\bibitem{ref23} Nwankwegu, A. S., Orji, M. U., & Onwosi, C. O. (2016). Studies on organic and in-organic biostimulants in bioremediation of diesel-contaminated arable soil. Chemosphere, 162, 148–156. https:// doi. org/ 10. 1016/j. chemo sphere. 2016. 07. 050

\bibitem{ref24} Nurfarahin, A. H., Mohamed, M. S., & Phang, L. Y. (2019). Development of palm fatty acid dis-tillate-containing medium for biosurfactant production by Pseudomonas sp. LM19. Molecules, 24(14), Article 2613. https:// doi. org/ 10. 3390/ molec ules2 41426 13

\bibitem{ref25} Morikawa, M., Hirata, Y., & Imanaka, T. (2000). A study on the structure–function relationship of lipopeptide biosurfactants. Biochimica et Biophysica Acta (BBA) - Molecular and Cell Biology of Lipids, 1488(3), 211–218. https:// doi. org/ 10. 1016/ S1388- 1981(00) 00117-2

\bibitem{ref26} Jain, D., Collins-Thompson, D., Lee, H., & Trevors, J. (1991). A drop-collapsing test for screen-ing surfactant-producing microorganisms. Journal of Microbiological Methods, 13(4), 271–279. https:// doi. org/ 10. 1016/ 0167- 7012(91) 90064-M

\bibitem{ref27} Morán, A. C., Martinez, M. A., & Siñeriz, F. (2002). Quantification of surfactin in culture super-natants by hemolytic activity. Biotechnology Letters, 24, 177–180. https:// doi. org/ 10. 1023/A: 10140 36408 116

\bibitem{ref28} Salimi, F., Khorshidi, M., Amirahmadi, F., & Amirahmadi, A. (2023). Effectiveness of phosphate and zinc solubilizing Paenarthrobacter nitroguajacolicus P1 as halotolerant rhizobacterium with growth-promoting activity on Pistacia vera L. Current Microbiology, 80(10), 336.

\bibitem{ref29} Abid, S., Farid, A., Abid, R., Rehman, M. U., Alsanie, W. F., Alhomrani, M., et al. (2022). Identifi-cation, biochemical characterization, and safety attributes of locally isolated Lactobacillus fermen-tum from Bubalus bubalis (buffalo) milk as a probiotic. Microorganisms, 10(5), Article 954. https:// doi. org/ 10. 3390/ micro organ isms1 00509 54

\bibitem{ref30} Lane, D. (1991). 16S/23S rRNA sequencing. In E. Stackebrandt, & M. Goodfellow (Eds.), Nucleic acid techniques in bacterial systematics (pp. 115–175). John Wiley and Sons.

\bibitem{ref31} Cooper, D. G., & Goldenberg, B. G. (1987). Surface-active agents from two Bacillus species. Applied and Environmental Microbiology, 53(2), 224–229. https:// doi. org/ 10. 1128/ aem. 53.2. 224- 229. 1987

\bibitem{ref32} Jain, R. M., Mody, K., Mishra, A., & Jha, B. (2012). Physicochemical characterization of biosur-factant and its potential to remove oil from soil and cotton cloth. Carbohydrate Polymers, 89(4), 1110–1116. https:// doi. org/ 10. 1016/j. carbp ol. 2012. 03. 071

\bibitem{ref33} Jain, R. M., Mody, K., Mishra, A., & Jha, B. (2012). Isolation and structural characterization of bio-surfactant produced by an alkaliphilic bacterium Cronobacter sakazakii isolated from oil contaminated wastewater. Carbohydrate Polymers, 87(3), 2320–2326. https:// doi. org/ 10. 1016/j. carbp ol. 2011. 10. 062

\bibitem{ref34} Lee, B.-B., Chan, E.-S., Ravindra, P., & Khan, T. A. (2012). Surface tension of viscous biopolymer solutions measured using the du Nouy ring method and the drop weight methods. Polymer Bulletin, 69, 471–489. https:// doi. org/ 10. 1007/ s00289- 011- 0719-x

\bibitem{ref35} Patowary, K., Patowary, R., Kalita, M. C., & Deka, S. (2017). Characterization of biosurfactant pro-duced during degradation of hydrocarbons using crude oil as sole source of carbon. Frontiers in Micro-biology, 8, Article 279. https:// doi. org/ 10. 3389/ fmicb. 2017. 00279

\bibitem{ref36} Zargar, A. N., Lymperatou, A., Skiadas, I., Kumar, M., & Srivastava, P. (2022). Structural and func-tional characterization of a novel biosurfactant from Bacillus sp. IITD106. Journal of Hazardous Materials, 423, Article 127201. https:// doi. org/ 10. 1016/j. jhazm at. 2021. 127201

\bibitem{ref37} Jui, A. H., Bhuiyan, M. N. I., Bhowmik, B., Khatun, N., Chowdhury, A., Bhuiyan, R. H., et al. (2024). Exploration and characterization of a newly isolated bacterium, Enterobacter quasihormaechei strain BDIFST24001, capable of producing rhamnolipid biosurfactant for oil remediation. Access Microbiol-ogy, 6(8), 000830, v4. https:// doi. org/ 10. 1099/ acmi.0. 000830

\bibitem{ref38} Thavasi, R., Sharma, S., & Jayalakshmi, S. (2011). Evaluation of screening methods for the isolation of biosurfactant producing marine bacteria. Journal of Petroleum & Environmental Biotechnology, 1(2), 1–7. https:// doi. org/ 10. 4172/ 2157- 7463. 10001 04

\bibitem{ref39} Wang, Q., Wang, J., Zhu, H., Zhuo, J., Liu, D., & Shen, Q. (2017). Purification and identification of the biosurfactant produced by Enterobacter hormaechei B4 and its application. Scientia Agricultura Sinica, 50(22), 4350–4361. https:// doi. org/ 10. 3864/j. issn. 0578- 1752. 2017. 22. 011

\bibitem{ref40} Muneeswari, R., Swathi, K., Sekaran, G., & Ramani, K. (2022). Microbial-induced biosurfactant-mediated biocatalytic approach for the bioremediation of simulated marine oil spill.  International Journal of Environmental Science and Technology, 1–14. https:// doi. org/ 10. 1007/ s13762- 022- 04214-4

\bibitem{ref41} Fardami, A., Kawo, A., Yahaya, S., Ahmad, S., & Ibrahim, U. (2022). Stability and optimization of biosurfactant production by Enterobacter cloacae AYF1. African Journal of Biotechnology and Applied Research, 1(1), 157–167.

\bibitem{ref42} Lovaglio, R. B., dos Santos, F. J., Junior, M. J., & Contiero, J. (2011). Rhamnolipid emulsifying activ-ity and emulsion stability: PH rules. Colloids and Surfaces B: Biointerfaces, 85(2), 301–305. https:// doi. org/ 10. 1016/j. colsu rfb. 2011. 03. 012

\bibitem{ref43} Mulligan, C. N. (2005). Environmental applications for biosurfactants. Environmental Pollution, 133(2), 183–198. https:// doi. org/ 10. 1016/j. envpol. 2004. 06. 009

\bibitem{ref44} Gudiña, E. J., Teixeira, J. A., & Rodrigues, L. R. (2016). Biosurfactants produced by marine microor-ganisms with therapeutic applications. Marine Drugs, 14(2), 38. https:// doi. org/ 10. 3390/ md140 20038

\bibitem{ref45} Nitschke, M., & Costa, S. (2007). Biosurfactants in food industry. Trends in Food Science & Technol- ogy, 18(5), 252–259. https:// doi. org/ 10. 1016/j. tifs. 2007. 02. 002

\bibitem{ref46} Rosen, M. J., & Kunjappu, J. T. (2012). Surfactants and interfacial phenomena. John Wiley & Sons.

\bibitem{ref47} Bai, G., Brusseau, M. L., & Miller, R. M. (1997). Biosurfactant-enhanced removal of residual hydro-carbon from soil. Journal of Contaminant Hydrology, 25(1–2), 157–170. https:// doi. org/ 10. 1016/ S0169- 7722(96) 00035-2

\bibitem{ref48} Banat, I. M., Franzetti, A., Gandolfi, I., Bestetti, G., Martinotti, M. G., Fracchia, L., et  al. (2010). Microbial biosurfactants production, applications and future potential. Applied Microbiology and Bio-technology, 87, 427–444. https:// doi. org/ 10. 1007/ s00253- 010- 2589-0

\bibitem{ref49} Elazzazy, A. M., Abdelmoneim, T., & Almaghrabi, O. (2015). Isolation and characterization of biosur-factant production under extreme environmental conditions by alkali-halo-thermophilic bacteria from 7430 Applied Biochemistry and Biotechnology (2025) 197:7409–7430 Saudi Arabia. Saudi Journal of Biological Sciences, 22(4), 466–475. https:// doi. org/ 10. 1016/j. sjbs. 2014. 11. 003

\bibitem{ref50} Banat, I. M., Franzetti, A., Gandolfi, I., Bestetti, G., Martinotti, M. G., Fracchia, L., et~al. (2010). Microbial biosurfactants production, applications and future potential. Applied Microbiology and Bio-technology, 87, 427–444. https:// doi. org/ 10. 1007/ s00253- 010- 2589-0

\bibitem{ref51} Beal, R., & Betts, W. (2000). Role of rhamnolipid biosurfactants in the uptake and mineralization of hexadecane in Pseudomonas aeruginosa. Journal of Applied Microbiology, 89(1), 158–168. https:// doi. org/ 10. 1046/j. 1365- 2672. 2000. 01105.x

\bibitem{ref52} Sheppard, J. D., & Mulligan, C. N. (1987). The production of surfactin by Bacillus subtilis grown on peat hydrolysate. Applied Microbiology and Biotechnology, 27, 110–116. https:// doi. org/ 10. 1007/ BF002 57795

\bibitem{ref53} Gudiña, E. J., Fernandes, E. C., Rodrigues, A. I., Teixeira, J. A., & Rodrigues, L. R. (2015). Biosur-factant production by Bacillus subtilis using corn steep liquor as culture medium. Frontiers in Microbi-ology, 6, 59. https:// doi. org/ 10. 3389/ fmicb. 2015. 00059

\bibitem{ref54} Khoshdast, H., Sam, A., Vali, H., & Noghabi, K. A. (2011). Effect of rhamnolipid biosurfactants on performance of coal and mineral flotation. International Biodeterioration & Biodegradation, 65(8), 1238–1243. https:// doi. org/ 10. 1016/j. ibiod. 2011. 10. 002

\bibitem{ref55} Soberón-Chávez, G., & Maier, R. M. (2011). Biosurfactants: A general overview. In Biosurfactants: From genes to applications (pp. 1–11). Springer. https:// doi. org/ 10. 1007/ 978-3- 642- 14490-5_1

\bibitem{ref56} Morais, I., Cordeiro, A., Teixeira, G., Domingues, V., Nardi, R., Monteiro, A., et~al. (2017). Biological and physicochemical properties of biosurfactants produced by Lactobacillus jensenii P 6A and Lacto-bacillus gasseri P 65. Microbial Cell Factories, 16, 1–15. https:// doi. org/ 10. 1186/ s12934- 017- 0640-1

\bibitem{ref57} Legawiec, K. J., Kruszelnicki, M., Zawadzka, M., Basařová, P., Zawala, J., & Polowczyk, I. (2023). Towards green flotation: Investigating the effect of rhamnolipid biosurfactant on single bubble adhe-sion dynamics. Journal Of Molecular Liquids, 388, Article 122759. https:// doi. org/ 10. 1016/j. molliq. 2023. 122759

\bibitem{ref58} Santos, A., Silva, M., Costa, E., Rufino, R., Santos, V., Ramos, C., et~al. (2018). Production and char-acterization of a biosurfactant produced by Streptomyces sp. DPUA 1559 isolated from lichens of the Amazon region. Brazilian Journal of Medical and Biological Research, 51(2), Article e6657. https:// doi. org/ 10. 1590/ 1414- 431x2 01766 57

\bibitem{ref59} Santos, A., Silva, M., Costa, E., Rufino, R., Santos, V., Ramos, C., et al. (2018). Production and char-acterization of a biosurfactant produced by Streptomyces sp. DPUA 1559 isolated from lichens of the Amazon region. Brazilian Journal of Medical and Biological Research, 51(2), Article e6657. https:// doi. org/ 10. 1590/ 1414- 431x2 01766 57

\bibitem{ref60} Oyibo, N. W. S., Akani, N. P., & Ugboma, C. J. (2020). Biodegradation of total petroleum hydrocarbon by molecularly identified bacteria isolated from an oilfield wastewater in Nigeria. Open Access Journal of Microbiology & Biotechnology, 5(2), 1–7. https:// doi. org/ 10. 23880/ OAJMB- 16000 164

\bibitem{ref61} Ammar, A. B., Bouassida, M., Bouallegue, A., Fourati, N., Gerardi, G., Muñiz, P., et al. (2023). Iso-lation and characterization of two glycolipopeptids biosurfactants produced by a Lactiplantibacillus plantarum OL5 strain isolated from green olive curing water. World Journal of Microbiology and Bio-technology, 39(11), 308. https:// doi. org/ 10. 1007/ s11274- 023- 03525-w

\bibitem{ref62} Pandey, R., Krishnamurthy, B., Singh, H. P., & Batish, D. R. (2022). Evaluation of a glycolipopeptide biosurfactant from Aeromonas hydrophila RP1 for bioremediation and enhanced oil recovery. Journal of Cleaner Production, 345, Article 131098. https:// doi. org/ 10. 1016/j. jclep ro. 2022. 131098

\bibitem{ref63} Sharma, V., Singh, D., Manzoor, M., Banpurkar, A. G., Satpute, S. K., & Sharma, D. (2022). Char-acterization and cytotoxicity assessment of biosurfactant derived from Lactobacillus pentosus NCIM 2912. Brazilian Journal of Microbiology, 53(1), 327–340. https:// doi. org/ 10. 1007/ s42770- 021- 00551-z

\bibitem{ref64} Gharaie, S., Ohadi, M., Hassanshahian, M., Shakibaie, M., Shahriary, P., & Forootanfar, H. (2023). Glycolipopeptide biosurfactant from Bacillus pumilus SG: Physicochemical characterization, optimi-zation, antibiofilm and antimicrobial activity evaluation.  3 Biotech, 13(10), 321.  https:// doi. org/ 10. 1007/ s13205- 023- 03554-5

\bibitem{ref65} Gil, C. V., Rebocho, A. T., Esmail, A., Sevrin, C., Grandfils, C., Torres, C. A., et  al. (2022). Char-acterization of the thermostable biosurfactant produced by Burkholderia thailandensis DSM 13276. Polymers, 14(10), 2088. https:// doi. org/ 10. 3390/ polym 14102 088

\end{thebibliography}
\end{document}
